% ============================================================================
\chapter{x86 架构速查}
\label{ch:appendix-x86-ref}
% ============================================================================

\section{寄存器}

\subsection{通用寄存器（32 位）}

\begin{longtable}{lll}
  \toprule
  \textbf{寄存器} & \textbf{惯用用途} & \textbf{调用约定} \\
  \midrule
  EAX & 累加器、返回值    & caller-saved \\
  EBX & 基址             & callee-saved \\
  ECX & 计数器           & caller-saved \\
  EDX & 数据、I/O        & caller-saved \\
  ESI & 源索引           & callee-saved \\
  EDI & 目的索引         & callee-saved \\
  EBP & 栈帧基指针       & callee-saved \\
  ESP & 栈指针           & 特殊 \\
  \bottomrule
\end{longtable}

\subsection{段寄存器}

\begin{longtable}{ll}
  \toprule
  \textbf{寄存器} & \textbf{用途} \\
  \midrule
  CS & 代码段选择子（隐含 CPL） \\
  DS & 默认数据段 \\
  SS & 栈段 \\
  ES & 附加数据段 \\
  FS / GS & 通用附加段（Linux 用 GS 做 per-CPU / TLS） \\
  \bottomrule
\end{longtable}

\subsection{控制寄存器}

\begin{longtable}{ll}
  \toprule
  \textbf{寄存器} & \textbf{关键位} \\
  \midrule
  CR0 & bit 0 (PE): 保护模式 \\
      & bit 31 (PG): 开启分页 \\
  CR2 & Page Fault 时的故障虚拟地址 \\
  CR3 & 页目录物理地址 \\
  CR4 & bit 4 (PSE): 4MiB 大页支持 \\
      & bit 7 (PGE): 全局页 \\
  \bottomrule
\end{longtable}

\section{GDT 条目格式}

% TODO: 8 字节条目的完整 bitfield 图

\section{IDT 条目格式}

% TODO: 8 字节中断门的完整 bitfield 图

\section{PDE / PTE 格式}

\begin{bitfield}
  \bit{0}{P (Present) — 1=有效，0=访问触发 \#PF}
  \bit{1}{R/W — 1=读写，0=只读}
  \bit{2}{U/S — 1=用户态可访问，0=仅内核}
  \bit{3}{PWT — 写透缓存策略}
  \bit{4}{PCD — 禁用缓存}
  \bit{5}{A (Accessed) — CPU 访问后自动置 1}
  \bit{6}{D (Dirty) — CPU 写入后自动置 1（仅 PTE）}
  \bit{7}{PS/PAT — PDE: 4MiB 大页; PTE: PAT}
  \bit{8}{G (Global) — TLB 不随 CR3 切换而刷新}
  \bit{9--11}{AVL — 软件可用位}
  \bit{12--31}{页帧物理地址高 20 位}
\end{bitfield}

\section{Page Fault 错误码}

\begin{bitfield}
  \bit{0}{P — 0=页不存在，1=权限违规}
  \bit{1}{W/R — 0=读操作，1=写操作}
  \bit{2}{U/S — 0=内核态，1=用户态}
  \bit{3}{RSVD — 页表条目保留位被置位}
  \bit{4}{I/D — 取指令引起}
\end{bitfield}

\section{中断向量分配}

\begin{longtable}{rll}
  \toprule
  \textbf{向量} & \textbf{名称} & \textbf{说明} \\
  \midrule
  0   & \#DE & 除零 \\
  3   & \#BP & 断点（int3） \\
  6   & \#UD & 无效操作码 \\
  8   & \#DF & Double Fault \\
  13  & \#GP & General Protection \\
  14  & \#PF & Page Fault \\
  32--47 & IRQ0--15 & 硬件中断（PIC remap 后） \\
  128 & --- & 系统调用（int 0x80） \\
  \bottomrule
\end{longtable}
